% =========================================================================
% Seguimiento de riesgos — 03-08-2026 (S2, segunda retrospectiva).
% Documento independiente de la memoria; creado copiando TEMPLATE.tex
% (ver README.md de esta carpeta). La retrospectiva de S2 tocaba el
% domingo 02-08 y se realizó el lunes 03-08 por indisponibilidad del
% autor durante el fin de semana (compromisos laborales); el desfase se
% registra aquí como observación.
% Compilar DESDE ESTA CARPETA (2 pases, sin biber):
%   pdflatex --enable-installer -interaction=nonstopmode -file-line-error
%     -output-directory=build 2026-08-03-S2.tex
% =========================================================================
\documentclass[11pt,a4paper]{article}

% =========================================================================
% Preámbulo COMPARTIDO de los documentos independientes de riesgos
% (docs/thesis/plan-riesgos/ y docs/thesis/seguimiento-riesgos/).
%
% Estos documentos NO forman parte de main.tex (la memoria): son entregables
% sueltos en clase article, con el diseño heredado de los .tex de Gestión de
% Proyectos del autor, pero adaptados a la configuración global de la
% memoria (Carlito, márgenes 2,5 cm, parskip 6 pt, títulos de la guía).
% Un solo preámbulo para toda la familia: un ajuste de estilo aquí se
% propaga a todos los seguimientos futuros (uno por sprint hasta S37).
%
% RUTAS: los \input son relativos a la carpeta del documento, no a esta;
% compilar SIEMPRE desde la carpeta del .tex (plan-riesgos/ o
% seguimiento-riesgos/):
%   pdflatex --enable-installer -interaction=nonstopmode -file-line-error
%     -output-directory=build <documento>.tex        (2 pases; sin biber)
% =========================================================================

% --- Codificación (idéntico a packages.tex de la memoria) ---
\usepackage[utf8]{inputenc}
\usepackage[T1]{fontenc}

% --- Idioma: solo español (aquí no hay resumen trilingüe) ---
% es-noshorthands por coherencia con la memoria: acentos en UTF-8, sin
% caracteres activos (el conflicto con biblatex-apa no aplica aquí, pero
% mantener una única configuración de babel evita sorpresas al copiar texto
% entre la memoria y estos documentos).
\usepackage[spanish, es-tabla, es-noshorthands]{babel}

% --- Márgenes: 2,5 cm laterales e inferior, como la memoria ---
% El margen superior se reparte para alojar la cabecera enmarcada del
% diseño de referencia (2 cm de caja + separación) sin comerse el cuerpo:
% el texto arranca a 4 cm del borde y el ancho de texto sigue siendo 16 cm.
\usepackage[a4paper, left=2.5cm, right=2.5cm, bottom=2.5cm,
            top=1.2cm, includehead, headheight=2.2cm, headsep=0.6cm]{geometry}

% --- Tipografía: Carlito, la misma decisión que en style.tex ---
% (clon métrico libre de Calibri; pdflatex no puede usar la real).
% No se replica el arreglo de versalitas de style.tex: aquí no hay
% bibliografía APA y ningún texto usa \textsc.
\usepackage[sfdefault]{carlito}

% --- Espaciado: interlineado sencillo + 6 pt entre párrafos (guía) ---
\usepackage[skip=6pt]{parskip}

% --- Tablas ---
% longtable: el registro y el plan de respuesta superan una página.
% AVISO ESPERADO: "ignored error: Infinite glue shrinkage" en cada salto de
% página de una longtable (bug conocido e inofensivo, latex3/latex2e#1907;
% ver el comentario largo en packages.tex de la memoria).
\usepackage{array}
\usepackage{longtable}
\usepackage{tabularx}
% \RaggedRight (ragged2e) en vez de \raggedright en columnas muy estrechas:
% el \raggedright de LaTeX desactiva de facto la partición silábica y una
% palabra larga («Identificado») desborda una columna de 1,2 cm; la
% variante de ragged2e mantiene el guionado.
\usepackage{ragged2e}
% Las columnas X se componen como m{...}: el contenido de cada celda queda
% centrado verticalmente (aspecto de formulario del diseño de referencia).
\renewcommand{\tabularxcolumn}[1]{m{#1}}
% Tipos de columna del diseño de referencia:
%   L/C = ancho fijo alineado izquierda/centro; Y/Z = X alineada izq./centro.
% En columnas estrechas se prefiere Y (raggedright) a X justificada: la
% justificación en 3-4 cm abre ríos de blancos.
\newcolumntype{L}[1]{>{\raggedright\arraybackslash}m{#1}}
\newcolumntype{C}[1]{>{\centering\arraybackslash}m{#1}}
\newcolumntype{Y}{>{\raggedright\arraybackslash}X}
\newcolumntype{Z}{>{\centering\arraybackslash}X}
% Filas con aire en todas las tablas (el diseño de referencia lo fijaba
% tabla a tabla; aquí una sola vez).
\renewcommand{\arraystretch}{1.25}

% --- Utilidades ---
\usepackage{xcolor}    % \markfield (resaltado) y \todonote (macros.tex)
\usepackage{amssymb}   % \square, \blacksquare, \circlearrowright
\usepackage{xspace}    % lo necesita \bukit (macros.tex)
\usepackage{enumitem}
\setlist[itemize]{leftmargin=*, topsep=2pt, itemsep=2pt}
\setlist[enumerate]{leftmargin=*, topsep=2pt, itemsep=2pt}

% --- Títulos según la guía de estilo, como en style.tex de la memoria ---
% (sin \chapter: la clase article no lo tiene; el resto, idéntico).
\usepackage{titlesec}
\titleformat{\section}
  {\normalfont\bfseries\fontsize{12pt}{14.4pt}\selectfont\raggedright}
  {\thesection.}{0.5em}{\MakeUppercase}
\titleformat{\subsection}
  {\normalfont\bfseries\fontsize{12pt}{14.4pt}\selectfont}
  {\thesubsection.}{0.5em}{}
\titleformat{\subsubsection}
  {\normalfont\itshape\fontsize{12pt}{14.4pt}\selectfont}
  {\thesubsubsection.}{0.5em}{}

% --- Paginación dinámica (regla estricta del autor, 26-07-2026) ---
% PROHIBIDO el salto de página manual también en estos documentos: el
% espacio lo garantizan estos ganchos (o un \Needspace* local ante un bloque
% indivisible), nunca un \newpage. Los .tex de referencia sí usaban
% \newpage; es la única desviación deliberada respecto a ellos.
\usepackage{needspace}
\AddToHook{cmd/section/before}{\Needspace*{5\baselineskip}}
\AddToHook{cmd/subsection/before}{\Needspace*{5\baselineskip}}
\AddToHook{cmd/subsubsection/before}{\Needspace*{4\baselineskip}}
\widowpenalty=10000
\clubpenalty=10000
\displaywidowpenalty=10000
\raggedbottom
% Estos formularios componen mucho texto en columnas de 2-4 cm: se concede
% a TeX un margen de estirado de emergencia antes de desbordar la línea
% (evita Overfull \hbox menores sin tocar el contenido).
\setlength{\emergencystretch}{1em}

% --- Datos globales y macros de la memoria (única fuente: no duplicar) ---
% =========================================================================
% Variables globales del documento.
% Cambiar un dato AQUÍ lo actualiza en todo el documento (portada,
% metadatos del PDF...). No escribir nunca estos datos "a mano" en otros
% archivos: usar siempre estos comandos.
% =========================================================================

% --- Trabajo ---
\newcommand{\thesisTitle}{BUKIT: red social orientada a objetivos}% TODO: título definitivo
\newcommand{\thesisSubtitle}{Memoria del Trabajo de Fin de Grado}%
\newcommand{\thesisDegree}{Grado en Ingeniería Informática de Gestión y Sistemas de Información}%
\newcommand{\thesisCourse}{2026--2027}% curso académico
\newcommand{\thesisDate}{Bilbao, \today}% TODO: <Lugar, día, mes, año> de la entrega/defensa
\newcommand{\thesisYear}{2027}% año de elaboración y defensa

% --- Personas (la plantilla del centro usa "Apellido1 Apellido2, Nombre") ---
\newcommand{\thesisAuthor}{Liviu Deleanu}% nombre natural (para metadatos del PDF)
\newcommand{\thesisAuthorFormal}{Deleanu, Liviu}% TODO: confirmar formato con dos apellidos
\newcommand{\thesisDirector}{Villamañe Gironés, Mikel}%
\newcommand{\thesisCoDirector}{Apellido1 Apellido2, Nombre}% TODO: solo si hay codirector/a

% --- Universidad ---
\newcommand{\thesisUniversity}{Universidad del País Vasco / Euskal Herriko Unibertsitatea}%
\newcommand{\thesisFaculty}{Escuela de Ingeniería de Bilbao}%

% --- Palabras clave del resumen trilingüe ---
\newcommand{\thesisKeywordsEs}{red social, objetivos, monorepo, TypeScript, PostgreSQL}% TODO
\newcommand{\thesisKeywordsEu}{sare soziala, helburuak, monorepo, TypeScript, PostgreSQL}% TODO: revisar euskera
\newcommand{\thesisKeywordsEn}{social network, goals, monorepo, TypeScript, PostgreSQL}% TODO

% =========================================================================
% Comandos propios. Definir aquí cualquier cosa que se repita en el texto:
% si mañana cambia cómo se escribe, se cambia en un único sitio.
% =========================================================================

% Nombre del producto, siempre con el mismo formato.
\newcommand{\bukit}{\textsc{Bukit}\xspace}

% Términos en inglés, en cursiva de forma consistente: \eng{backend}
\newcommand{\eng}[1]{\textit{#1}}

% Notas de trabajo, visibles en rojo en el PDF. Buscar "todonote" y
% eliminarlas todas antes de entregar.
\newcommand{\todonote}[1]{\textcolor{red}{\textbf{[TODO: #1]}}}


% --- Cabecera y pie (diseño de referencia, sin logos) ---
% La caja de cabecera replica la del .tex de referencia (celda ancha de
% contexto + celda estrecha de identidad). La celda derecha llevaba el logo
% del centro: aquí va texto porque los logos UPV/EHU siguen pendientes de
% la fase de diseño (ver style.tex); cuando existan, sustituir por
% \includegraphics. 11 + 4,9 = 15,9 cm < 16 cm de texto: sin desbordes.
\usepackage{fancyhdr}
\newcommand{\riesgosHeader}{%
  \setlength{\tabcolsep}{0pt}%
  \renewcommand{\arraystretch}{1}%
  \begin{tabular}{|p{11cm}|p{4.9cm}|}
    \hline
    \parbox[c][2cm][c]{11cm}{\centering\bfseries
      {\fontsize{16}{19}\selectfont TRABAJO DE FIN DE GRADO\\[4pt]Curso \thesisCourse}%
    } &
    \parbox[c][2cm][c]{4.9cm}{\centering
      {\fontsize{14}{17}\selectfont\bfseries\bukit}\\[3pt]
      {\small\thesisFaculty}%
    }\\
    \hline
  \end{tabular}%
}
\pagestyle{fancy}
\fancyhf{}
\renewcommand{\headrulewidth}{0pt}
\fancyhead[C]{\riesgosHeader}
\fancyfoot[C]{\thepage}
% La primera página usa estilo "plain": se redefine para que también lleve
% la cabecera (mismo truco que el .tex de referencia).
\fancypagestyle{plain}{%
  \fancyhf{}%
  \renewcommand{\headrulewidth}{0pt}%
  \fancyhead[C]{\riesgosHeader}%
  \fancyfoot[C]{\thepage}%
}

% --- Símbolos y campos de formulario del diseño de referencia ---
\newcommand{\fieldline}[1]{\underline{\hspace{#1}}}
\newcommand{\bigline}{\noindent\rule{\textwidth}{0.4pt}}
\newcommand{\checkbox}{\(\square\)}
\newcommand{\checkedbox}{\(\blacksquare\)}
% Símbolo ⟳ = "si se rellena este apartado, nueva versión del plan".
% El original lo resolvía con un SVG externo (paquete svg + Inkscape);
% aquí basta el símbolo de amssymb: cero dependencias nuevas.
\newcommand{\updatesymbol}{\raisebox{-0.2ex}{\(\circlearrowright\)}}
% Resaltado amarillo del identificador de riesgo en los títulos.
\newcommand{\markfield}[1]{\colorbox{yellow!35}{\strut #1}}

% --- Enlaces y metadatos del PDF (al final, como exige hyperref) ---
% El pdftitle lo fija cada documento con \hypersetup tras cargar esto.
\usepackage[hidelinks]{hyperref}
\hypersetup{pdfauthor={\thesisAuthor}}

\hypersetup{pdftitle={Seguimiento de riesgos — 03-08-2026 (S2)}}

\begin{document}

\begin{center}
  {\LARGE\bfseries Seguimiento de riesgos — S2}\\[6pt]
\end{center}

\textbf{Fecha:} 03-08-2026 · \textbf{Sprint:} S2 (27-07 \(\to\) 02-08-2026; retrospectiva realizada el lunes 03-08 por compromisos laborales del fin de semana) · \textbf{Responsable:} \thesisAuthor

Regla de revisión: \textbf{críticos} cada sprint; \textbf{significativos} cada 2 sprints; \textbf{moderados} cada 4 sprints; cualquier riesgo, \textbf{de inmediato}, si se observa su indicador clave.

\textbf{El símbolo \updatesymbol\ indica que, si se rellena ese apartado, hay que emitir una nueva versión completa del plan de riesgos} en \texttt{../plan-riesgos/} y registrar el cambio en la tabla de seguimiento del cap.~7 de la memoria.

Riesgos revisados: \textbf{R1–R7 por cadencia} (críticos, semanal) y \textbf{R8–R11 por cadencia} (significativos, cada dos sprints: primera revisión por cadencia de R8, R9 y R11; R10 repite tras su revisión por disparador en S1).
Los moderados (R12, R13) tocan en el cierre de S4.
Datos del sprint: velocidad de S2 = 3 SP (US-02, cerrada el 27-07); horas reales cronometradas por el autor = 5~h~40~min de las 10~h declaradas (57~\%), con viernes--domingo sin actividad por compromisos laborales.

\bigline

\section{Medidas de mitigación (prevención)}

Para cada medida de los riesgos que tocan revisión en este seguimiento: ¿se está aplicando?
Si no, justificar y decidir si se mantiene/modifica/elimina.

% La tabla se divide en dos (críticos / significativos) porque tabularx no
% parte entre páginas y las 11 filas de S2 juntas desbordan la página
% (Overfull \vbox); la división coincide además con los dos grupos de
% cadencia revisados en este seguimiento.
% Tablas partidas en bloques de 3-4 filas (no \newpage, regla de paginación
% dinámica): tabularx es una caja indivisible y con las celdas de evidencia
% de S2 los 7 críticos juntos superan la altura de una página entera
% (Overfull \vbox); cada bloque repite la cabecera y salta con \Needspace.
\Needspace*{24\baselineskip}
\underline{\textbf{Riesgos críticos (R1--R7, cadencia semanal)}}

\begin{tabularx}{\textwidth}{|C{0.9cm}|Y|Y|C{2.2cm}|L{2.4cm}|}
\hline
\multicolumn{1}{|C{0.9cm}|}{\textbf{Id}} &
\multicolumn{1}{>{\centering\arraybackslash}X|}{\textbf{Medida de mitigación}} &
\multicolumn{1}{>{\centering\arraybackslash}X|}{\textbf{¿Aplicada?}} &
\multicolumn{1}{C{2.2cm}|}{\textbf{Justificación (del NO)}} &
\multicolumn{1}{C{2.4cm}|}{\textbf{Decisión}} \\
\hline
R1 & Toda idea nueva entra por el backlog. & \checkedbox\ Sí (EP17, publicidad simulada, entró el 28-07 por el backlog como opcional, con su RF-20 y su objetivo secundario en el cap.~3; el sprint no ganó puntos).\newline \checkbox\ No & — & \checkedbox\ Mantener\newline \checkbox\ Modificar\newline \checkbox\ Eliminar \\
\hline
R2 & Medir velocidad real desde el primer sprint y recalibrar en cada cierre. & \checkedbox\ Sí, parcial: la medición sí (segunda: 3 SP en S2); la recalibración de horas/punto, aún no.\newline \checkbox\ No & Dos datos no marcan tendencia (el indicador de R2 exige 3 sprints); se recalibra al cierre de S3. & \checkedbox\ Mantener\newline \checkbox\ Modificar\newline \checkbox\ Eliminar \\
\hline
R3 & Diario al cierre de cada sesión; memoria con hueco en todos los sprints. & \checkedbox\ Sí (diario el mismo día en las cuatro jornadas de S2; la memoria recibió la revisión completa del 28-07 —62 correcciones— y la lectura manual del autor hasta el cap.~8; el cap.~8 documenta la US-02).\newline \checkbox\ No & — & \checkedbox\ Mantener\newline \checkbox\ Modificar\newline \checkbox\ Eliminar \\
\hline
R4 & Citar al escribir; pendientes con nota visible. & \checkedbox\ Sí (la revisión del 28-07 citó en el momento todo lo detectado, con 8 entradas nuevas de bibliografía; los pendientes viven como \texttt{todonote}).\newline \checkbox\ No & — & \checkedbox\ Mantener\newline \checkbox\ Modificar\newline \checkbox\ Eliminar \\
\hline
\end{tabularx}

\Needspace*{20\baselineskip}
\underline{\textbf{Riesgos críticos (continuación)}}

\begin{tabularx}{\textwidth}{|C{0.9cm}|Y|Y|C{2.2cm}|L{2.4cm}|}
\hline
\multicolumn{1}{|C{0.9cm}|}{\textbf{Id}} &
\multicolumn{1}{>{\centering\arraybackslash}X|}{\textbf{Medida de mitigación}} &
\multicolumn{1}{>{\centering\arraybackslash}X|}{\textbf{¿Aplicada?}} &
\multicolumn{1}{C{2.2cm}|}{\textbf{Justificación (del NO)}} &
\multicolumn{1}{C{2.4cm}|}{\textbf{Decisión}} \\
\hline
R5 & Revisar todo lo generado antes de integrarlo; nada cerrado sin comprobaciones. & \checkedbox\ Sí (US-02 pasó revisión adversarial antes del commit —9 defectos corregidos— y CI en verde; el autor validó además la memoria con lectura manual propia).\newline \checkbox\ No & — & \checkedbox\ Mantener\newline \checkbox\ Modificar\newline \checkbox\ Eliminar \\
\hline
R6 & Investigación previa acotada antes de tecnología desconocida. & \checkedbox\ Sí (pino/pino-http era terreno nuevo acotado y la US-02 cerró dentro de su estimación).\newline \checkbox\ No & — & \checkedbox\ Mantener\newline \checkbox\ Modificar\newline \checkbox\ Eliminar \\
\hline
R7 & Minimización desde el diseño (RGPD). & \checkedbox\ Sí (sin modelo de datos aún; además, la regla del KYC quedó precisada por escrito: la imagen del documento no se almacena nunca en el sistema).\newline \checkbox\ No & — & \checkedbox\ Mantener\newline \checkbox\ Modificar\newline \checkbox\ Eliminar \\
\hline
\end{tabularx}

\Needspace*{22\baselineskip}
\underline{\textbf{Riesgos significativos (R8--R11, cadencia quincenal)}}

\begin{tabularx}{\textwidth}{|C{0.9cm}|Y|Y|C{2.2cm}|L{2.4cm}|}
\hline
\multicolumn{1}{|C{0.9cm}|}{\textbf{Id}} &
\multicolumn{1}{>{\centering\arraybackslash}X|}{\textbf{Medida de mitigación}} &
\multicolumn{1}{>{\centering\arraybackslash}X|}{\textbf{¿Aplicada?}} &
\multicolumn{1}{C{2.2cm}|}{\textbf{Justificación (del NO)}} &
\multicolumn{1}{C{2.4cm}|}{\textbf{Decisión}} \\
\hline
R8 & Horas reales por sesión en el diario; cada retrospectiva las compara con las declaradas. & \checkedbox\ Sí (cronometraje directo del autor por sesión; esta retrospectiva compara: 5~h~40~min de 10~h declaradas).\newline \checkbox\ No & — & \checkedbox\ Mantener\newline \checkbox\ Modificar\newline \checkbox\ Eliminar \\
\hline
R9 & Integraciones externas como objetivos secundarios o contra versiones simuladas; proveedores con capa gratuita. & \checkedbox\ Sí (sin integraciones aún; el principio de versiones simuladas, fijado en el cap.~8 al cierre de S1, quedó validado durante S2 por la revisión del 28-07 y la lectura manual del autor).\newline \checkbox\ No & — & \checkedbox\ Mantener\newline \checkbox\ Modificar\newline \checkbox\ Eliminar \\
\hline
R10 & Margen temporal deliberado + documentación continua que permite retomar sin pérdida. & \checkedbox\ Sí (retomar hoy costó leer la última entrada del diario; el margen sigue bastando: el parón de 3 días no consume margen adicional al de la semana ya absorbida en S1).\newline \checkbox\ No & — & \checkedbox\ Mantener\newline \checkbox\ Modificar\newline \checkbox\ Eliminar \\
\hline
R11 & Todo en git con remoto y subidas frecuentes. & \checkedbox\ Sí (las cuatro jornadas de S2 acabaron commiteadas y subidas —Días 5 a 8; la sesión del 28-07, partida, subió la mañana del 29 como Día 6—; árbol limpio al abrir esta retrospectiva).\newline \checkbox\ No & — & \checkedbox\ Mantener\newline \checkbox\ Modificar\newline \checkbox\ Eliminar \\
\hline
\end{tabularx}

\Needspace*{6\baselineskip}
\underline{\textbf{Cambios realizados en las medidas de mitigación}} \updatesymbol

\begin{tabularx}{\textwidth}{|C{1.4cm}|Y|Y|Y|}
\hline
\multicolumn{1}{|C{1.4cm}|}{\textbf{Id}} &
\multicolumn{1}{>{\centering\arraybackslash}X|}{\textbf{Medida original}} &
\multicolumn{1}{>{\centering\arraybackslash}X|}{\textbf{Cambio realizado}} &
\multicolumn{1}{>{\centering\arraybackslash}X|}{\textbf{Justificación}} \\
\hline
— & — & — & — \\
\hline
\end{tabularx}

\Needspace*{6\baselineskip}
\underline{\textbf{Nuevas medidas de mitigación a añadir}} \updatesymbol

\begin{tabularx}{\textwidth}{|C{1.4cm}|Y|Y|}
\hline
\multicolumn{1}{|C{1.4cm}|}{\textbf{Id}} &
\multicolumn{1}{>{\centering\arraybackslash}X|}{\textbf{Nueva medida}} &
\multicolumn{1}{>{\centering\arraybackslash}X|}{\textbf{Justificación}} \\
\hline
— & — & — \\
\hline
\end{tabularx}

\bigline

\section{Riesgos a revisar en este seguimiento}

% ------------------ R1 ------------------
\Needspace*{12\baselineskip}
\subsection*{Riesgo: \markfield{R1 — Crecimiento descontrolado del alcance}}

\textbf{Indicadores (clave y de resolución)}

\begin{tabularx}{\textwidth}{|Z|Z|Z|Z|Z|}
\hline
\textbf{Indicador} & \textbf{Valor esperado} & \textbf{Valor real} & \textbf{¿Cumplido?} & \textbf{Acción / observaciones} \\
\hline
Sprint gana >2 SP tras la planificación / función fuera del backlog. & 0 & 0 (EP17 entró por el backlog como opcional; los refinamientos de reglas de la lectura manual son precisiones, no funciones nuevas). & \checkbox\ Sí\quad \checkedbox\ No & Sin acción. \\
\hline
\end{tabularx}

Si el \textbf{indicador clave} se cumple: reevaluar probabilidad/impacto y severidad.
Si el \textbf{indicador de resolución} se cumple: cerrar el riesgo.

\textbf{Nueva probabilidad (1–3):} \fieldline{0.8cm} · \textbf{Nuevo impacto (1–3):} \fieldline{0.8cm} · \textbf{Nueva severidad:} \checkbox\ Crítico · \checkbox\ Significativo · \checkbox\ Moderado\ \updatesymbol

\textbf{¿Cerrar el riesgo?} \checkbox\ Sí · \checkedbox\ No\ \updatesymbol

\textbf{¿Se ha materializado?} \checkbox\ Sí · \checkedbox\ No\ \updatesymbol

\textbf{Nuevo estado} (cap.~7): \checkedbox\ Sin novedad · \checkbox\ Activo · \checkbox\ Materializado · \checkbox\ Resuelto · \checkbox\ Cerrado

\textbf{Plan de contingencia (solo si se ha materializado)}

\begin{tabularx}{\textwidth}{|Y|Y|Y|}
\hline
\multicolumn{1}{|>{\centering\arraybackslash}X|}{\textbf{Medida de contingencia}} &
\multicolumn{1}{>{\centering\arraybackslash}X|}{\textbf{¿Aplicada?}} &
\multicolumn{1}{>{\centering\arraybackslash}X|}{\textbf{Justificación (del NO)}} \\
\hline
 & \checkbox\ Sí\quad \checkbox\ No & \\
\hline
\end{tabularx}

Reflexión sobre el funcionamiento de la contingencia: no aplicable.

% ------------------ R2 ------------------
\Needspace*{12\baselineskip}
\subsection*{Riesgo: \markfield{R2 — Error sistemático de estimación del esfuerzo}}

\textbf{Indicadores (clave y de resolución)}

\begin{tabularx}{\textwidth}{|Z|Z|Z|Z|Z|}
\hline
\textbf{Indicador} & \textbf{Valor esperado} & \textbf{Valor real} & \textbf{¿Cumplido?} & \textbf{Acción / observaciones} \\
\hline
Velocidad de los 3 primeros sprints desviada >25~\% de la hipótesis. & — (necesita 3 sprints). & Segunda medición: 3 SP en S2. La jornada que cerró la US-02 sumó 1~h~20 en total (planning, migración de pruebas y revisión de los caps.~1--2 incluidos): la historia consumió como mucho esa cota, \(\leq\)0,45~h/SP —y sensiblemente menos en realidad— frente a la hipótesis de 2,7. Dos mediciones consecutivas muy por debajo. & \checkbox\ Sí\quad \checkedbox\ No & El indicador se evalúa formalmente al cierre de S3; con dos datos coincidentes, anticipar que la equivalencia horas/punto se recalibrará a la baja (la asistencia de IA acelera la ejecución muy por encima de lo previsto). \\
\hline
\end{tabularx}

Si el \textbf{indicador clave} se cumple: reevaluar probabilidad/impacto y severidad.
Si el \textbf{indicador de resolución} se cumple: cerrar el riesgo.

\textbf{Nueva probabilidad (1–3):} \fieldline{0.8cm} · \textbf{Nuevo impacto (1–3):} \fieldline{0.8cm} · \textbf{Nueva severidad:} \checkbox\ Crítico · \checkbox\ Significativo · \checkbox\ Moderado\ \updatesymbol

\textbf{¿Cerrar el riesgo?} \checkbox\ Sí · \checkedbox\ No\ \updatesymbol

\textbf{¿Se ha materializado?} \checkbox\ Sí · \checkedbox\ No\ \updatesymbol

\textbf{Nuevo estado} (cap.~7): \checkedbox\ Sin novedad · \checkbox\ Activo · \checkbox\ Materializado · \checkbox\ Resuelto · \checkbox\ Cerrado

\textbf{Plan de contingencia (solo si se ha materializado)}

\begin{tabularx}{\textwidth}{|Y|Y|Y|}
\hline
\multicolumn{1}{|>{\centering\arraybackslash}X|}{\textbf{Medida de contingencia}} &
\multicolumn{1}{>{\centering\arraybackslash}X|}{\textbf{¿Aplicada?}} &
\multicolumn{1}{>{\centering\arraybackslash}X|}{\textbf{Justificación (del NO)}} \\
\hline
 & \checkbox\ Sí\quad \checkbox\ No & \\
\hline
\end{tabularx}

Reflexión sobre el funcionamiento de la contingencia: no aplicable.

% ------------------ R3 ------------------
\Needspace*{12\baselineskip}
\subsection*{Riesgo: \markfield{R3 — Memoria relegada frente al desarrollo}}

\textbf{Indicadores (clave y de resolución)}

\begin{tabularx}{\textwidth}{|Z|Z|Z|Z|Z|}
\hline
\textbf{Indicador} & \textbf{Valor esperado} & \textbf{Valor real} & \textbf{¿Cumplido?} & \textbf{Acción / observaciones} \\
\hline
Capítulo afectado 2 sprints sin actualizar / diario >48~h tras la sesión. & 0 & 0 (diario el mismo día en las cuatro jornadas; en S2 la memoria fue el foco: revisión completa + lectura manual). & \checkbox\ Sí\quad \checkedbox\ No & Observación: la retrospectiva se desplazó del domingo al lunes por indisponibilidad; el indicador mide el diario de las sesiones, que no se vio afectado. \\
\hline
\end{tabularx}

Si el \textbf{indicador clave} se cumple: reevaluar probabilidad/impacto y severidad.
Si el \textbf{indicador de resolución} se cumple: cerrar el riesgo.

\textbf{Nueva probabilidad (1–3):} \fieldline{0.8cm} · \textbf{Nuevo impacto (1–3):} \fieldline{0.8cm} · \textbf{Nueva severidad:} \checkbox\ Crítico · \checkbox\ Significativo · \checkbox\ Moderado\ \updatesymbol

\textbf{¿Cerrar el riesgo?} \checkbox\ Sí · \checkedbox\ No\ \updatesymbol

\textbf{¿Se ha materializado?} \checkbox\ Sí · \checkedbox\ No\ \updatesymbol

\textbf{Nuevo estado} (cap.~7): \checkedbox\ Sin novedad · \checkbox\ Activo · \checkbox\ Materializado · \checkbox\ Resuelto · \checkbox\ Cerrado

\textbf{Plan de contingencia (solo si se ha materializado)}

\begin{tabularx}{\textwidth}{|Y|Y|Y|}
\hline
\multicolumn{1}{|>{\centering\arraybackslash}X|}{\textbf{Medida de contingencia}} &
\multicolumn{1}{>{\centering\arraybackslash}X|}{\textbf{¿Aplicada?}} &
\multicolumn{1}{>{\centering\arraybackslash}X|}{\textbf{Justificación (del NO)}} \\
\hline
 & \checkbox\ Sí\quad \checkbox\ No & \\
\hline
\end{tabularx}

Reflexión sobre el funcionamiento de la contingencia: no aplicable.

% ------------------ R4 ------------------
\Needspace*{12\baselineskip}
\subsection*{Riesgo: \markfield{R4 — Incumplimiento de la normativa de la memoria}}

\textbf{Indicadores (clave y de resolución)}

\begin{tabularx}{\textwidth}{|Z|Z|Z|Z|Z|}
\hline
\textbf{Indicador} & \textbf{Valor esperado} & \textbf{Valor real} & \textbf{¿Cumplido?} & \textbf{Acción / observaciones} \\
\hline
Afirmación no original sin cita detectada en revisión. & 0 & La revisión completa del 28-07 encontró afirmaciones sin cita —todas corregidas en esa misma sesión, con 8 entradas nuevas de bibliografía—, de modo que el literal vigente del indicador se cumplió; la lectura manual del autor (preliminares a cap.~8) no encontró ninguna adicional. & \checkedbox\ Sí\quad \checkbox\ No & Reevaluación abajo, sin cambios; y el indicador se redefine \updatesymbol: una detección corregida al momento es la revisión funcionando, no la señal del riesgo. \\
\hline
\end{tabularx}

Si el \textbf{indicador clave} se cumple: reevaluar probabilidad/impacto y severidad.
Si el \textbf{indicador de resolución} se cumple: cerrar el riesgo.

\textbf{Nueva probabilidad (1–3):} \underline{\ 2\ } · \textbf{Nuevo impacto (1–3):} \underline{\ 3\ } · \textbf{Nueva severidad:} \checkedbox\ Crítico · \checkbox\ Significativo · \checkbox\ Moderado\ \updatesymbol\ (sin cambio: la reevaluación confirma la valoración P2\(\times\)I3)

\textbf{¿Cerrar el riesgo?} \checkbox\ Sí · \checkedbox\ No\ \updatesymbol

\textbf{¿Se ha materializado?} \checkbox\ Sí · \checkedbox\ No\ \updatesymbol

Cambio de \textbf{plan} (no de estado) \updatesymbol: el indicador clave se redefine — de «una revisión de capítulo encuentra una afirmación no original sin su cita» a «\ldots y esta no queda corregida en la misma sesión que la detecta» — porque en S1 y S2 el literal saltó por detecciones corregidas al momento, con lo que medía el funcionamiento de la revisión y no la señal del riesgo.
Recogido en el plan v2 (emitido hoy) y en la tabla de seguimiento del cap.~7.

\textbf{Nuevo estado} (cap.~7): \checkedbox\ Sin novedad (el cambio es del plan, no del estado) · \checkbox\ Activo · \checkbox\ Materializado · \checkbox\ Resuelto · \checkbox\ Cerrado

\textbf{Plan de contingencia (solo si se ha materializado)}

\begin{tabularx}{\textwidth}{|Y|Y|Y|}
\hline
\multicolumn{1}{|>{\centering\arraybackslash}X|}{\textbf{Medida de contingencia}} &
\multicolumn{1}{>{\centering\arraybackslash}X|}{\textbf{¿Aplicada?}} &
\multicolumn{1}{>{\centering\arraybackslash}X|}{\textbf{Justificación (del NO)}} \\
\hline
 & \checkbox\ Sí\quad \checkbox\ No & \\
\hline
\end{tabularx}

Reflexión sobre el funcionamiento de la contingencia: no aplicable.

% ------------------ R5 ------------------
\Needspace*{12\baselineskip}
\subsection*{Riesgo: \markfield{R5 — Dependencia acrítica de la IA generativa}}

\textbf{Indicadores (clave y de resolución)}

\begin{tabularx}{\textwidth}{|Z|Z|Z|Z|Z|}
\hline
\textbf{Indicador} & \textbf{Valor esperado} & \textbf{Valor real} & \textbf{¿Cumplido?} & \textbf{Acción / observaciones} \\
\hline
Decisión/bloque integrado que no puede explicarse sin la herramienta. & 0 & 0 (la US-02 quedó razonada en el diario con sus alternativas; la lectura manual de la memoria por el autor y su lista propia de fallos son evidencia directa de comprensión y criterio sobre lo generado). & \checkbox\ Sí\quad \checkedbox\ No & Refuerzo metodológico de S2: los fallos de la lectura manual se aplican solo tras validar el entendimiento mutuo con un plan explicado. \\
\hline
\end{tabularx}

Si el \textbf{indicador clave} se cumple: reevaluar probabilidad/impacto y severidad.
Si el \textbf{indicador de resolución} se cumple: cerrar el riesgo.

\textbf{Nueva probabilidad (1–3):} \fieldline{0.8cm} · \textbf{Nuevo impacto (1–3):} \fieldline{0.8cm} · \textbf{Nueva severidad:} \checkbox\ Crítico · \checkbox\ Significativo · \checkbox\ Moderado\ \updatesymbol

\textbf{¿Cerrar el riesgo?} \checkbox\ Sí · \checkedbox\ No\ \updatesymbol

\textbf{¿Se ha materializado?} \checkbox\ Sí · \checkedbox\ No\ \updatesymbol

\textbf{Nuevo estado} (cap.~7): \checkedbox\ Sin novedad · \checkbox\ Activo · \checkbox\ Materializado · \checkbox\ Resuelto · \checkbox\ Cerrado

\textbf{Plan de contingencia (solo si se ha materializado)}

\begin{tabularx}{\textwidth}{|Y|Y|Y|}
\hline
\multicolumn{1}{|>{\centering\arraybackslash}X|}{\textbf{Medida de contingencia}} &
\multicolumn{1}{>{\centering\arraybackslash}X|}{\textbf{¿Aplicada?}} &
\multicolumn{1}{>{\centering\arraybackslash}X|}{\textbf{Justificación (del NO)}} \\
\hline
 & \checkbox\ Sí\quad \checkbox\ No & \\
\hline
\end{tabularx}

Reflexión sobre el funcionamiento de la contingencia: no aplicable.

% ------------------ R6 ------------------
\Needspace*{12\baselineskip}
\subsection*{Riesgo: \markfield{R6 — Curva de aprendizaje de la pila tecnológica}}

\textbf{Indicadores (clave y de resolución)}

\begin{tabularx}{\textwidth}{|Z|Z|Z|Z|Z|}
\hline
\textbf{Indicador} & \textbf{Valor esperado} & \textbf{Valor real} & \textbf{¿Cumplido?} & \textbf{Acción / observaciones} \\
\hline
Historia\slash investigación al doble de su tiempo previsto. & 0 & 0 (pino/pino-http era terreno nuevo y la US-02 cerró dentro de una jornada de 1~h~20 en total, muy por debajo de su estimación). & \checkbox\ Sí\quad \checkedbox\ No & Sin acción. \\
\hline
\end{tabularx}

Si el \textbf{indicador clave} se cumple: reevaluar probabilidad/impacto y severidad.
Si el \textbf{indicador de resolución} se cumple: cerrar el riesgo.

\textbf{Nueva probabilidad (1–3):} \fieldline{0.8cm} · \textbf{Nuevo impacto (1–3):} \fieldline{0.8cm} · \textbf{Nueva severidad:} \checkbox\ Crítico · \checkbox\ Significativo · \checkbox\ Moderado\ \updatesymbol

\textbf{¿Cerrar el riesgo?} \checkbox\ Sí · \checkedbox\ No\ \updatesymbol

\textbf{¿Se ha materializado?} \checkbox\ Sí · \checkedbox\ No\ \updatesymbol

\textbf{Nuevo estado} (cap.~7): \checkedbox\ Sin novedad · \checkbox\ Activo · \checkbox\ Materializado · \checkbox\ Resuelto · \checkbox\ Cerrado

\textbf{Plan de contingencia (solo si se ha materializado)}

\begin{tabularx}{\textwidth}{|Y|Y|Y|}
\hline
\multicolumn{1}{|>{\centering\arraybackslash}X|}{\textbf{Medida de contingencia}} &
\multicolumn{1}{>{\centering\arraybackslash}X|}{\textbf{¿Aplicada?}} &
\multicolumn{1}{>{\centering\arraybackslash}X|}{\textbf{Justificación (del NO)}} \\
\hline
 & \checkbox\ Sí\quad \checkbox\ No & \\
\hline
\end{tabularx}

Reflexión sobre el funcionamiento de la contingencia: no aplicable.

% ------------------ R7 ------------------
\Needspace*{12\baselineskip}
\subsection*{Riesgo: \markfield{R7 — Incumplimiento del RGPD en el diseño}}

\textbf{Indicadores (clave y de resolución)}

\begin{tabularx}{\textwidth}{|Z|Z|Z|Z|Z|}
\hline
\textbf{Indicador} & \textbf{Valor esperado} & \textbf{Valor real} & \textbf{¿Cumplido?} & \textbf{Acción / observaciones} \\
\hline
Dato personal sin finalidad declarada en el modelo / requisito que exija almacenar el DNI. & 0 & 0 (el modelo de datos aún no existe; en S2 la regla de minimización del KYC quedó además precisada por escrito: la imagen del documento no se almacena nunca). & \checkbox\ Sí\quad \checkedbox\ No & Primera revisión real cuando se diseñe la entidad de usuario (US-04). \\
\hline
\end{tabularx}

Si el \textbf{indicador clave} se cumple: reevaluar probabilidad/impacto y severidad.
Si el \textbf{indicador de resolución} se cumple: cerrar el riesgo.

\textbf{Nueva probabilidad (1–3):} \fieldline{0.8cm} · \textbf{Nuevo impacto (1–3):} \fieldline{0.8cm} · \textbf{Nueva severidad:} \checkbox\ Crítico · \checkbox\ Significativo · \checkbox\ Moderado\ \updatesymbol

\textbf{¿Cerrar el riesgo?} \checkbox\ Sí · \checkedbox\ No\ \updatesymbol

\textbf{¿Se ha materializado?} \checkbox\ Sí · \checkedbox\ No\ \updatesymbol

\textbf{Nuevo estado} (cap.~7): \checkedbox\ Sin novedad · \checkbox\ Activo · \checkbox\ Materializado · \checkbox\ Resuelto · \checkbox\ Cerrado

\textbf{Plan de contingencia (solo si se ha materializado)}

\begin{tabularx}{\textwidth}{|Y|Y|Y|}
\hline
\multicolumn{1}{|>{\centering\arraybackslash}X|}{\textbf{Medida de contingencia}} &
\multicolumn{1}{>{\centering\arraybackslash}X|}{\textbf{¿Aplicada?}} &
\multicolumn{1}{>{\centering\arraybackslash}X|}{\textbf{Justificación (del NO)}} \\
\hline
 & \checkbox\ Sí\quad \checkbox\ No & \\
\hline
\end{tabularx}

Reflexión sobre el funcionamiento de la contingencia: no aplicable.

% ------------------ R8 ------------------
\Needspace*{12\baselineskip}
\subsection*{Riesgo: \markfield{R8 — Capacidad real inferior a la declarada}}

\textbf{Indicadores (clave y de resolución)}

\begin{tabularx}{\textwidth}{|Z|Z|Z|Z|Z|}
\hline
\textbf{Indicador} & \textbf{Valor esperado} & \textbf{Valor real} & \textbf{¿Cumplido?} & \textbf{Acción / observaciones} \\
\hline
Horas reales <70~\% de las declaradas dos semanas consecutivas sin causa puntual identificada. & 0 & Dos semanas <70~\% (S1: 2~h~40, 27~\%; S2: 5~h~40, 57~\%), pero ambas \textbf{con causa puntual identificada} (S1: materialización de R10; S2: compromisos laborales de viernes a domingo). & \checkbox\ Sí\quad \checkedbox\ No & Vigilancia expresa para S3: si el déficit se repite por trabajo, la causa dejará de considerarse puntual (patrón estructural del verano), el indicador se dará por cumplido y tocará aplicar la contingencia (rebajar la capacidad declarada de F1 al nivel observado). \\
\hline
\end{tabularx}

Si el \textbf{indicador clave} se cumple: reevaluar probabilidad/impacto y severidad.
Si el \textbf{indicador de resolución} se cumple: cerrar el riesgo.

\textbf{Nueva probabilidad (1–3):} \fieldline{0.8cm} · \textbf{Nuevo impacto (1–3):} \fieldline{0.8cm} · \textbf{Nueva severidad:} \checkbox\ Crítico · \checkbox\ Significativo · \checkbox\ Moderado\ \updatesymbol

\textbf{¿Cerrar el riesgo?} \checkbox\ Sí · \checkedbox\ No\ \updatesymbol

\textbf{¿Se ha materializado?} \checkbox\ Sí · \checkedbox\ No\ \updatesymbol

\textbf{Nuevo estado} (cap.~7): \checkedbox\ Sin novedad · \checkbox\ Activo · \checkbox\ Materializado · \checkbox\ Resuelto · \checkbox\ Cerrado

\textbf{Plan de contingencia (solo si se ha materializado)}

\begin{tabularx}{\textwidth}{|Y|Y|Y|}
\hline
\multicolumn{1}{|>{\centering\arraybackslash}X|}{\textbf{Medida de contingencia}} &
\multicolumn{1}{>{\centering\arraybackslash}X|}{\textbf{¿Aplicada?}} &
\multicolumn{1}{>{\centering\arraybackslash}X|}{\textbf{Justificación (del NO)}} \\
\hline
 & \checkbox\ Sí\quad \checkbox\ No & \\
\hline
\end{tabularx}

Reflexión sobre el funcionamiento de la contingencia: no aplicable.

% ------------------ R9 ------------------
\Needspace*{12\baselineskip}
\subsection*{Riesgo: \markfield{R9 — Dependencia de servicios externos}}

\textbf{Indicadores (clave y de resolución)}

\begin{tabularx}{\textwidth}{|Z|Z|Z|Z|Z|}
\hline
\textbf{Indicador} & \textbf{Valor esperado} & \textbf{Valor real} & \textbf{¿Cumplido?} & \textbf{Acción / observaciones} \\
\hline
Cambio de precios/límites de un servicio, agotamiento de capa gratuita o indisponibilidad >24~h. & 0 & 0 (sin servicios externos integrados todavía; sin anuncios de cambios en los candidatos). & \checkbox\ Sí\quad \checkedbox\ No & Primera revisión con datos reales al integrar el primer servicio (US-03: destino de despliegue). \\
\hline
\end{tabularx}

Si el \textbf{indicador clave} se cumple: reevaluar probabilidad/impacto y severidad.
Si el \textbf{indicador de resolución} se cumple: cerrar el riesgo.

\textbf{Nueva probabilidad (1–3):} \fieldline{0.8cm} · \textbf{Nuevo impacto (1–3):} \fieldline{0.8cm} · \textbf{Nueva severidad:} \checkbox\ Crítico · \checkbox\ Significativo · \checkbox\ Moderado\ \updatesymbol

\textbf{¿Cerrar el riesgo?} \checkbox\ Sí · \checkedbox\ No\ \updatesymbol

\textbf{¿Se ha materializado?} \checkbox\ Sí · \checkedbox\ No\ \updatesymbol

\textbf{Nuevo estado} (cap.~7): \checkedbox\ Sin novedad · \checkbox\ Activo · \checkbox\ Materializado · \checkbox\ Resuelto · \checkbox\ Cerrado

\textbf{Plan de contingencia (solo si se ha materializado)}

\begin{tabularx}{\textwidth}{|Y|Y|Y|}
\hline
\multicolumn{1}{|>{\centering\arraybackslash}X|}{\textbf{Medida de contingencia}} &
\multicolumn{1}{>{\centering\arraybackslash}X|}{\textbf{¿Aplicada?}} &
\multicolumn{1}{>{\centering\arraybackslash}X|}{\textbf{Justificación (del NO)}} \\
\hline
 & \checkbox\ Sí\quad \checkbox\ No & \\
\hline
\end{tabularx}

Reflexión sobre el funcionamiento de la contingencia: no aplicable.

% ------------------ R10 ------------------
\Needspace*{12\baselineskip}
\subsection*{Riesgo: \markfield{R10 — Indisponibilidad prolongada del autor}}

\textbf{Indicadores (clave y de resolución)}

\begin{tabularx}{\textwidth}{|Z|Z|Z|Z|Z|}
\hline
\textbf{Indicador} & \textbf{Valor esperado} & \textbf{Valor real} & \textbf{¿Cumplido?} & \textbf{Acción / observaciones} \\
\hline
>3 días consecutivos sin poder trabajar por causa sobrevenida. & 0 & 3 días justos (31-07 \(\to\) 02-08, compromisos laborales): en el umbral, sin superarlo. & \checkbox\ Sí\quad \checkedbox\ No & El episodio de S1 sigue superado (Resuelto). Consecuencia práctica del parón: la retrospectiva se desplazó al lunes 03-08. Si el próximo parón supera los 3 días, revisión inmediata. \\
\hline
\end{tabularx}

Si el \textbf{indicador clave} se cumple: reevaluar probabilidad/impacto y severidad.
Si el \textbf{indicador de resolución} se cumple: cerrar el riesgo.

\textbf{Nueva probabilidad (1–3):} \fieldline{0.8cm} · \textbf{Nuevo impacto (1–3):} \fieldline{0.8cm} · \textbf{Nueva severidad:} \checkbox\ Crítico · \checkbox\ Significativo · \checkbox\ Moderado\ \updatesymbol

\textbf{¿Cerrar el riesgo?} \checkbox\ Sí · \checkedbox\ No\ \updatesymbol

\textbf{¿Se ha materializado?} \checkbox\ Sí · \checkedbox\ No\ \updatesymbol

\textbf{Nuevo estado} (cap.~7): \checkedbox\ Sin novedad (el episodio de S1 permanece Resuelto) · \checkbox\ Activo · \checkbox\ Materializado · \checkbox\ Resuelto · \checkbox\ Cerrado

\textbf{Plan de contingencia (solo si se ha materializado)}

\begin{tabularx}{\textwidth}{|Y|Y|Y|}
\hline
\multicolumn{1}{|>{\centering\arraybackslash}X|}{\textbf{Medida de contingencia}} &
\multicolumn{1}{>{\centering\arraybackslash}X|}{\textbf{¿Aplicada?}} &
\multicolumn{1}{>{\centering\arraybackslash}X|}{\textbf{Justificación (del NO)}} \\
\hline
 & \checkbox\ Sí\quad \checkbox\ No & \\
\hline
\end{tabularx}

Reflexión sobre el funcionamiento de la contingencia: no aplicable.

% ------------------ R11 ------------------
\Needspace*{12\baselineskip}
\subsection*{Riesgo: \markfield{R11 — Pérdida de trabajo o del entorno local}}

\textbf{Indicadores (clave y de resolución)}

\begin{tabularx}{\textwidth}{|Z|Z|Z|Z|Z|}
\hline
\textbf{Indicador} & \textbf{Valor esperado} & \textbf{Valor real} & \textbf{¿Cumplido?} & \textbf{Acción / observaciones} \\
\hline
Anomalías del equipo / >2 sesiones consecutivas sin subir al remoto. & 0 & 0 (las cuatro jornadas de S2 —Días 5 a 8— acabaron commiteadas y subidas; la del 28-07, partida, subió la mañana del 29; racha máxima sin subir: una sesión; sin anomalías del equipo). & \checkbox\ Sí\quad \checkedbox\ No & Sin acción. \\
\hline
\end{tabularx}

Si el \textbf{indicador clave} se cumple: reevaluar probabilidad/impacto y severidad.
Si el \textbf{indicador de resolución} se cumple: cerrar el riesgo.

\textbf{Nueva probabilidad (1–3):} \fieldline{0.8cm} · \textbf{Nuevo impacto (1–3):} \fieldline{0.8cm} · \textbf{Nueva severidad:} \checkbox\ Crítico · \checkbox\ Significativo · \checkbox\ Moderado\ \updatesymbol

\textbf{¿Cerrar el riesgo?} \checkbox\ Sí · \checkedbox\ No\ \updatesymbol

\textbf{¿Se ha materializado?} \checkbox\ Sí · \checkedbox\ No\ \updatesymbol

\textbf{Nuevo estado} (cap.~7): \checkedbox\ Sin novedad · \checkbox\ Activo · \checkbox\ Materializado · \checkbox\ Resuelto · \checkbox\ Cerrado

\textbf{Plan de contingencia (solo si se ha materializado)}

\begin{tabularx}{\textwidth}{|Y|Y|Y|}
\hline
\multicolumn{1}{|>{\centering\arraybackslash}X|}{\textbf{Medida de contingencia}} &
\multicolumn{1}{>{\centering\arraybackslash}X|}{\textbf{¿Aplicada?}} &
\multicolumn{1}{>{\centering\arraybackslash}X|}{\textbf{Justificación (del NO)}} \\
\hline
 & \checkbox\ Sí\quad \checkbox\ No & \\
\hline
\end{tabularx}

Reflexión sobre el funcionamiento de la contingencia: no aplicable.

\bigline

\section{Reflexión desde el último seguimiento}

\textbf{¿Ha ocurrido algo reseñable desde el último seguimiento?} \checkedbox\ Sí · \checkbox\ No

S2 fue la semana de mayor actividad documental del proyecto: la US-02 se cerró el primer día del sprint con revisión adversarial previa al commit, la memoria pasó una revisión completa de coherencia (62 correcciones y 8 entradas nuevas de bibliografía), los documentos de riesgos migraron a LaTeX (plan v1 + este mecanismo de seguimiento), CLAUDE.md se realineó con la memoria y el autor realizó su lectura manual de preliminares a cap.~8, con una lista de fallos que se corrigen con un plan validado.
Velocidad de S2: 3 SP; horas: 5~h~40~min de 10 declaradas (57~\%), con viernes--domingo sin actividad por compromisos laborales y la retrospectiva desplazada al lunes 03-08.

\Needspace*{6\baselineskip}
\textbf{Nuevos riesgos identificados} \updatesymbol

\begin{tabularx}{\textwidth}{|C{1.4cm}|Y|C{0.8cm}|C{0.8cm}|C{1cm}|C{2.4cm}|}
\hline
\multicolumn{1}{|C{1.4cm}|}{\textbf{Id}} &
\multicolumn{1}{>{\centering\arraybackslash}X|}{\textbf{Descripción}} &
\multicolumn{1}{C{0.8cm}|}{\textbf{P}} &
\multicolumn{1}{C{0.8cm}|}{\textbf{I}} &
\multicolumn{1}{C{1cm}|}{\textbf{Exp.}} &
\multicolumn{1}{C{2.4cm}|}{\textbf{Severidad}} \\
\hline
— & — & — & — & — & — \\
\hline
\end{tabularx}

\textbf{¿Afecta algo a los objetivos o al alcance (cap.~3 de la memoria)?} \checkbox\ Sí · \checkedbox\ No

\textbf{Justificación:} las correcciones pendientes de la lectura manual son refinamientos de reglas de producto y de redacción, no cambios de objetivos ni de alcance; la consulta al tutor sobre la estructura del cap.~3 es documental; la posible re-base de la fecha de cierre (16/05/2027) se registrará cuando ocurra, con su versión del plan si procede.

\textbf{Otros eventos relevantes:} redefinición del indicador clave de R4 decidida en esta misma revisión (\updatesymbol, recogida en el plan v2 emitido hoy y en la tabla de seguimiento del cap.~7); bloqueo de \texttt{graphify.exe} por una directiva de control de aplicaciones de Windows (detectado al preparar esta retrospectiva; pendiente de desbloquear, sin pérdida de sesión que active R13, cuya revisión por cadencia toca en S4); el límite de uso del plan de suscripción de la herramienta de IA partió la sesión del 28-07 (fricción asumida y prevista por el plan contratado).

\end{document}
